% Chapter 2, Topic _Linear Algebra_ Jim Hefferon
%  http://joshua.smcvt.edu/linearalgebra
%  2001-Jun-11
\topic{Lapangan}
\index{lapangan|(}
Perhitungan yang hanya melibatkan bilangan bulat atau hanya bilangan rasional
jauh lebih mudah daripada perhitungan yang melibatkan bilangan real.
Dapatkah struktur aljabar lain, seperti bilangan bulat atau bilangan rasional,
menggantikan
\( \Re \) dalam definisi ruang vektor?

Jika yang dimaksud dengan ``dapat'' adalah bahwa hasil-hasil dalam bab ini
tetap berlaku,
maka terdapat daftar syarat yang secara wajar harus dimiliki suatu struktur
(yakni suatu sistem bilangan)
agar dapat menggantikan $\Re$.
Sebuah \definend{lapangan}\index{lapangan!definisi} adalah himpunan
\( \F \) yang dilengkapi dengan operasi
`\( + \)'
dan `\( \cdot \)' sedemikian sehingga
\begin{tfae}
   \item untuk sembarang \( a,b\in\F \), hasil \( a+b \) berada di \( \F \), dan
         \( a+b=b+a \), dan
         jika \( c\in\F \), maka \( a+(b+c)=(a+b)+c \)
   \item untuk sembarang \( a,b\in\F \), hasil \( a\cdot b \) berada di \( \F \), dan
         \( a\cdot b=b\cdot a \), dan
         jika \( c\in\F \), maka \( a\cdot (b\cdot c)=(a\cdot b)\cdot c \)
   \item jika \( a,b,c\in\F \), maka \( a\cdot (b+c)=a\cdot b+a\cdot c \)
   \item terdapat unsur \( 0\in\F \) sedemikian sehingga
         jika \( a\in\F \), maka \( a+0=a \), dan
         untuk setiap \( a\in\F \) terdapat unsur \( -a\in\F \)
                sedemikian sehingga \( (-a)+a=0 \)
   \item terdapat unsur \( 1\in\F \) sedemikian sehingga
         jika \( a\in\F \), maka \( a\cdot 1=a \), dan
         untuk setiap unsur \( a\neq 0 \) di \( \F \)
                terdapat unsur \( a^{-1}\in\F \)
                sedemikian sehingga \( a^{-1}\cdot a=1 \).
\end{tfae}

Sebagai contoh, struktur aljabar yang terdiri atas himpunan bilangan real
beserta operasi penjumlahan dan perkalian yang lazim merupakan sebuah lapangan.
Contoh lapangan lainnya adalah himpunan bilangan rasional dengan operasi
penjumlahan dan perkalian yang lazim.
Contoh struktur aljabar yang bukan lapangan adalah
bilangan bulat, karena struktur itu tidak memenuhi syarat terakhir.

Beberapa contoh lainnya lebih mengejutkan.
Himpunan \( \mathbb{B}=\set{0,1} \) dengan operasi berikut:
\par\vfill\pagebreak
\begin{center}
  \begin{tabular}{c|cc}
    \( + \) &\( 0 \) &\( 1 \) \\
    \hline
    \( 0 \) &\( 0 \) &\( 1 \) \\
    \( 1 \) &\( 1 \) &\( 0 \)
  \end{tabular}
  \qquad
  \begin{tabular}{c|cc}
    \( \cdot \) &\( 0 \) &\( 1 \) \\
     \hline
       \( 0 \)  &\( 0 \) &\( 0 \)  \\
       \( 1 \)  &\( 0 \) &\( 1 \)
  \end{tabular}
\end{center}
merupakan sebuah lapangan; lihat \nearbyexercise{exer:BinField}.

Dalam buku ini kita dapat mengembangkan Aljabar Linear sebagai teori
ruang vektor dengan skalar dari sembarang lapangan.
Dalam hal itu,
hampir semua pernyataan di sini akan tetap berlaku jika
`\( \Re \)' diganti dengan `\( \F \)', yakni dengan
mengambil koefisien, entri vektor,
dan entri matriks sebagai unsur \( \F \)
(pengecualiannya adalah pernyataan yang melibatkan jarak atau sudut,
yang memerlukan pengembangan tambahan).
Berikut beberapa contohnya; masing-masing berlaku untuk ruang vektor \( V \)
di atas lapangan \( \F \).
\begin{itemize}
  %\makebox[\parindent]{\ \hfil\ }
  \item[$*$] Untuk sembarang \( \vec{v}\in V \) dan \( a\in\F \),
  (i)~\( 0\cdot\vec{v}=\zero \),
  (ii)~\( -1\cdot\vec{v}+\vec{v}=\zero \),
  dan (iii)~\( a\cdot\vec{0}=\vec{0} \).

  \item[$*$] Rentang, yaitu himpunan kombinasi linear, dari suatu himpunan
  bagian \( V \) merupakan subruang dari \( V \).

  \item[$*$] Setiap himpunan bagian dari suatu himpunan bebas linear juga
     bebas linear.

  \item[$*$] Dalam ruang vektor berdimensi hingga,
    sembarang dua basis memiliki jumlah unsur yang sama.
\end{itemize}
(Bahkan pernyataan yang tidak menyebutkan \( \F \) secara eksplisit
menggunakan sifat-sifat lapangan dalam pembuktiannya.)

Kita tidak akan mengembangkan ruang vektor dalam tatanan yang lebih umum ini
karena abstraksi tambahan tersebut dapat mengalihkan perhatian.
Gagasan yang ingin kita tonjolkan sudah tampak ketika kita membatasi diri
pada bilangan real.

Pengecualiannya adalah Bab Lima.
Di sana kita harus memfaktorkan polinom,
sehingga kita akan beralih untuk membahas ruang vektor di atas
lapangan bilangan kompleks.

\begin{exercises}
  \item 
    Periksalah bahwa bilangan real membentuk suatu lapangan.
    \begin{answer}
      Dengan menelaah kelima syarat tersebut, terlihat bahwa semuanya sudah
      dikenal dari matematika dasar.
    \end{answer}
  \item 
    Buktikan bahwa berikut ini merupakan lapangan.
    \begin{exparts*}
       \partsitem Bilangan rasional $\Q$
       \partsitem Bilangan kompleks  $\C$
    \end{exparts*}
    \begin{answer}
      Seperti pada soal sebelumnya, dengan menelaah kelima syarat tersebut
      terlihat bahwa, untuk kedua struktur ini,
      semua sifatnya sudah dikenal.
    \end{answer}
  \item 
     Berikan contoh yang menunjukkan bahwa sistem bilangan bulat
     bukan lapangan.
     \begin{answer}
       Bilangan bulat tidak memenuhi syarat~(5).
       Sebagai contoh, $2$ tidak memiliki invers perkalian\Dash meskipun
       $2$ adalah bilangan bulat, $1/2$ bukan bilangan bulat.
     \end{answer}
  \item \label{exer:BinField} 
     Periksalah bahwa himpunan $\mathbb{B}=\set{0,1}$ merupakan lapangan
     dengan operasi yang tercantum di atas.
     \begin{answer}
       Pemeriksaan ini dapat kita lakukan dengan mendaftarkan semua
       kemungkinan.
       Sebagai contoh, untuk memverifikasi bagian pertama syarat~(1), kita
       harus memeriksa bahwa struktur tersebut tertutup terhadap penjumlahan
       dan bahwa penjumlahan bersifat komutatif, $a+b=b+a$.
       Kita dapat memeriksa keduanya untuk semua pasangan $a$ dan~$b$ yang
       mungkin karena hanya ada empat pasangan seperti itu.
       Demikian pula, untuk asosiativitas hanya ada delapan tripel $a$, $b$,
       $c$, sehingga pemeriksaannya tidak terlalu panjang.
       (Ada cara lain untuk melakukan pemeriksaan tersebut; khususnya, Anda
       mungkin mengenali operasi ini sebagai aritmetika modulo~$2$.
       Namun, pemeriksaan menyeluruh tidaklah memberatkan.)
     \end{answer}
  \item 
     Berikan operasi yang sesuai agar himpunan $\set{0,1,2}$
     menjadi suatu lapangan.
     \begin{answer}
       Operasi berikut memenuhi syarat.
       \begin{center}
          \begin{tabular}{c|ccc}
             \( + \) &\( 0 \) &\( 1 \) &$2$ \\
             \hline
             \( 0 \) &\( 0 \) &\( 1 \) &$2$ \\
             \( 1 \) &\( 1 \) &\( 2 \) &$0$ \\
             $2$     &$2$     &$0$     &$1$
          \end{tabular}
          \qquad
          \begin{tabular}{c|ccc}
            \( \cdot \) &\( 0 \) &\( 1 \) &$2$ \\
            \hline
            \( 0 \)  &\( 0 \) &\( 0 \) &$0$ \\
            \( 1 \)  &\( 0 \) &\( 1 \) &$2$ \\
            $2$      &$0$     &$2$     &$1$
          \end{tabular}
       \end{center}
       Seperti pada soal sebelumnya, kita dapat memverifikasi bahwa operasi
       tersebut memenuhi semua syarat dengan mendaftarkan seluruh kasus.
     \end{answer}
\end{exercises}
\index{lapangan|)}
%\include{la}
%\end{document}
